% Part: incompleteness
% Chapter: representability-in-q
% Section: representing-relations

\documentclass[../../../include/open-logic-section]{subfiles}

\begin{document}

\olfileid{inc}{req}{rel}

\olsection{संबंधांचे निरूपण}

एखादा \emph{संबंध} निरूपणीय असणे म्हणजे काय, ते आता सांगू.

\begin{defn}
\ollabel{defn:representing-relations} नैसर्गिक संख्यांवरील
$R(x_0,\dots,x_k)$ हा संबंध {\em $\Th{Q}$ मध्ये निरूपणीय}
असतो, जर असे सूत्र $!A_R(x_0,\dots,x_k)$ असेल की
$R(n_0,\dots,n_k)$ सत्य असेल तेव्हा $\Th{Q}$ हे
$!A_R(\num{n_0},\dots,\num{n_k})$ सिद्ध करते, आणि
$R(n_0,\dots,n_k)$ असत्य असेल तेव्हा $\Th{Q}$ हे
$\lnot !A_R(\num{n_0},
\dots, \num{n_k})$ सिद्ध करते.
\end{defn}

\begin{thm}
\ollabel{thm:representing-rels} संबंध~$\Th{Q}$ मध्ये निरूपणीय
असतो तर आणि तरच तो संगणनक्षम असतो.
\end{thm}

\begin{proof}
प्रथम डावीकडून उजवीकडे जाणारी दिशा पाहू. समजा
$R(x_0,\dots,x_k)$ चे निरूपण $!A_R(x_0,\dots,x_k)$ या
सूत्राने होते. $R$ चे संगणन करण्यासाठी पुढील कार्यपद्धती
आहे: $n_0$, \dots,~$n_k$ ही आदाने मिळाल्यावर
$!A_R(\num{n_0}, \dots, \num{n_k})$ चा पुरावा आणि
$\lnot !A_R(\num{n_0}, \dots, \num{n_k})$ चा पुरावा
यांचा एकाच वेळी शोध घ्या. गृहीतकानुसार यांपैकी एक पुरावा
नक्की सापडेल. पहिला सापडल्यास ``होय'' आणि दुसरा सापडल्यास
``नाही'' असे उत्तर द्या. ही उपपत्ती सुसंगत असल्याने दोन्ही
विरुद्ध पुरावे मिळू शकत नाहीत.

उलट दिशेसाठी $R(x_0, \dots, x_k)$ संगणनक्षम आहे असे समजा.
व्याख्येनुसार याचा अर्थ $\Char{R}(x_0, \dots, x_k)$ हे
फलन संगणनक्षम आहे. \olref[int]{thm:representable-iff-comp} नुसार
$\Char{R}$ चे एखाद्या सूत्राने, समजा
$!A_{\Char{R}}(x_0, \dots, x_k,
y)$ ने, निरूपण होते. $!A_R(x_0, \dots, x_k)$ हे
$!A_{\Char{R}}(x_0,
\dots, x_k, \num{1})$ सूत्र असू द्या. मग कोणत्याही
$n_0$, \dots,~$n_k$ साठी $R(n_0,
\dots, n_k)$ सत्य असेल, तर $\Char{R}(n_0, \dots, n_k) = 1$.
त्या वेळी $\Th{Q}$ हे
$!A_{\Char{R}}(\num{n_0}, \dots, \num{n_k},
\num{1})$ सिद्ध करते; म्हणून $\Th{Q}$ हे
$!A_R(\num{n_0}, \dots,
\num{n_k})$ सिद्ध करते. दुसरीकडे $R(n_0, \dots, n_k)$
असत्य असेल, तर $\Char{R}(n_0, \dots, n_k) = 0$.
म्हणून $\Th{Q}$ पुढील विधान सिद्ध करते:
\[
\lforall[y][(!A_{\Char{R}}(\num{n_0}, \dots, \num{n_k}, y) \lif y =
  \num{0})].
\]
$\Th{Q}$ हे $\eq/[\num{0}][\num{1}]$ सिद्ध करत असल्याने,
$\Th{Q}$ हे $\lnot !A_{\Char{R}}(\num{n_0}, \dots, \num{n_k}, \num{1})$
सिद्ध करते; म्हणून $\lnot !A_R(\num{n_0}, \dots, \num{n_k})$
हेही सिद्ध करते.
\end{proof}

\begin{prob}
$R$ हा संबंध~$\Th{Q}$ मध्ये निरूपणीय असेल, तर
$\Char{R}$ हेसुद्धा निरूपणीय असते, हे दाखवा.
\end{prob}

\end{document}
